\documentclass[11pt,a4paper]{article}

\usepackage[T1]{fontenc}
\usepackage[utf8]{inputenc}
\usepackage[main=italian,provide=*]{babel}
\usepackage[a4paper,margin=2.2cm]{geometry}
\usepackage{booktabs}
\usepackage{array}
\usepackage{enumitem}
\usepackage{xcolor}
\usepackage{colortbl}

\definecolor{introcol}{RGB}{230,240,250}
\definecolor{riscol}{RGB}{234,246,234}

\pagestyle{empty}

\newcolumntype{L}[1]{>{\raggedright\arraybackslash}p{#1}}

\begin{document}

\begin{center}
{\Large\bfseries Selezione delle figure per slide}\\[4pt]
{\normalsize Proposta --- da confermare prima di produrre le slide}
\end{center}

\vspace{10pt}

\rowcolors{2}{white}{white}
\begin{center}
\begin{tabular}{c L{2.6cm} L{6.3cm} L{4.5cm}}
\toprule
\# & Slide & Figura proposta & Note \\
\midrule
\rowcolor{introcol!40}
1 & Frontespizio & --- & Nessuna figura, solo testo \\
\rowcolor{introcol!40}
2 & Contesto & --- & Nessun candidato esistente adatto (sarebbe uno schema concettuale nuovo, non una selezione) \\
\rowcolor{introcol!40}
3 & Stato dell'arte & --- & Nessun candidato --- eventuale riferimento testuale a Crippa et al. \\
\rowcolor{introcol!40}
4 & Obiettivi & --- & Slide di testo, non serve figura \\
\rowcolor{introcol!40}
5 & Modello fisico & \emph{(nessun candidato geometrico esistente)} & Cercato uno schema dimero/trimero (siti e legami): non esiste già fatto. Userei solo la formula dell'Hamiltoniana \\
\rowcolor{introcol!40}
6 & Metodo VQE & \texttt{circ\_ha.pdf} + \texttt{circ\_pma\_esteso.pdf} (dimero), affiancati & Verificati entrambi --- puliti, confronto diretto generico/simmetrico \\
\midrule
\rowcolor{riscol!40}
7 & Confronto ansatz & \texttt{fig03\_vqe\_ha\_vs\_pma.pdf} (dimero) & Verificata, messaggio chiarissimo \\
\rowcolor{riscol!40}
8 & Dinamica & \texttt{fig07\_correlatore\_validazione.pdf} (dimero) & Verificata, asse $t$ con unità esplicite \\
\rowcolor{riscol!40}
9 & Pipeline di rumore & \texttt{fig\_passo3\_F\_vs\_N.pdf} (dimero rumoroso) & Verificata --- compromesso e $N^*$ in un solo grafico \\
\rowcolor{riscol!40}
10 & Estensioni & \texttt{fig\_F\_vs\_N\_confronto.pdf} (dimero rumoroso) & Verificata --- ideale vs noise-aware, differenza trascurabile ben visibile \\
\rowcolor{riscol!40}
11 & Dimero vs trimero & --- & Nessuna figura esistente --- probabile tabella di confronto (già presente in \texttt{punti\_di\_lavoro\_riepilogo.tex}) \\
\rowcolor{riscol!40}
12 & Conclusioni & --- & Nessuna figura, solo testo \\
\bottomrule
\end{tabular}
\end{center}

\vspace{14pt}
\noindent\textbf{Punti aperti, da decidere:}
\begin{itemize}[nosep]
\item Slide 2--5: restano testuali, oppure si crea uno schema geometrico
      semplice nuovo (es. due pallini collegati per il dimero, un
      triangolo per il trimero) per la slide 5?
\item Slide 11: tabella di confronto testuale, o si costruisce una
      figura riassuntiva ad hoc?
\end{itemize}

\end{document}
